\subsection{Arbitrary smooth collars: an exact equation and the shape dependence}
The question on 67/021 says ``collar'' without prescribing its shape. We therefore treat any smooth embedded collar whose removal preserves the topological type, not only a collar of constant normal width. Curvature below is signed against the outward normal of the retained surface.

\begin{theorem}[Proved: the metric relation for every smooth collar]\label{A2-collar-general}
Let $(M,g)$ be a compact connected hyperbolic surface with geodesic boundary and $\chi(M)<0$. Let $M'\subset M$ be obtained by removing finitely many disjoint smooth open boundary collars, with $M'$ diffeomorphic to $M$. Write $k_g=g(\nabla_TT,n)$ on $\partial M'$, where $T$ is a unit tangent and $n$ is the outward unit normal. The unique conformally compatible hyperbolic metric with geodesic boundary is $g'=e^{2u}g$, where $u\in C^\infty(M')$ is the unique solution
\[
 \Delta_g u=e^{2u}-1,\qquad \partial_nu=k_g.
\]
Its area and boundary lengths are exactly
\[
 \int_{M'}e^{2u}\,dA_g=-2\pi\chi(M),\qquad
 L_{g'}(B)=\int_B e^u\,ds_g\quad(B\subset\partial M').
\]
If $k_g\ge0$ everywhere and is positive somewhere, then $u>0$ everywhere. In general the sign of $u$ is not determined by the fact that a collar was removed.
\end{theorem}
\begin{proof}
Uniformization of the bordered conformal surface supplies $g'$, smooth to the smooth boundary. The curvature formulas are
\[
 K_{e^{2u}g}=e^{-2u}(-1-\Delta_g u),\qquad
 k_{e^{2u}g}=e^{-u}(k_g-\partial_nu).
\]
Thus its defining conditions give the stated boundary-value problem. Conversely these formulas show that every smooth solution supplies the required metric. If $u_1,u_2$ are solutions, their difference $v$ satisfies $\Delta_gv=c(x)v$, $\partial_nv=0$, with
$c(x)=2\int_0^1e^{2((1-t)u_2+tu_1)}dt>0$. Multiplying by $v$ and integrating gives
$-\int|\nabla v|^2=\int cv^2$, so $v=0$.
Gauss--Bonnet for $g'$ gives its area, and conformal scaling gives the boundary-length formula. The maximum-principle and boundary-point argument in Theorem~\ref{A2-collar} uses only $k_g\ge0$, and proves the asserted strict positivity. The final assertion is proved by the following explicit family; no positivity is assumed when applying the displayed equation to a general collar.
\end{proof}

\begin{proposition}[Proved: first variation and a family refuting general monotonicity]\label{A2-collar-shape}
Let the boundary components of $(M,g)$ have arclength coordinates $s$, and choose smooth functions $a_i(s)>0$. In inward Fermi coordinates remove the collars
$0\le r<\varepsilon a_i(s)$, for sufficiently small $\varepsilon>0$, and let $u_\varepsilon$ be as in Theorem~\ref{A2-collar-general}. After any smooth identification with $M$ which is the identity at $\varepsilon=0$,
\[
 u_\varepsilon=\varepsilon h+O(\varepsilon^2)\quad\hbox{in }C^2,
 \qquad
 (\Delta_g-2)h=0,\quad \partial_nh=a_i-a_i''.
\]
Here the prime is differentiation in boundary arclength. In particular, take the unique hyperbolic pant all of whose boundary lengths are $2\pi$. On one boundary set
\[
 a_1(s)=1+\tfrac12\cos(ns),
\]
and on the other two set $a_2=a_3=1$. For every sufficiently large integer $n$, and then every sufficiently small $\varepsilon>0$, the resulting $u_\varepsilon$ takes a negative value at a point of the new first boundary. Thus even smooth analytic, everywhere inward collar removal need not increase the hyperbolic metric pointwise. The same metric takes a positive value somewhere.
\end{proposition}
\begin{proof}
In the original metric the collar is
$dr^2+\cosh^2r\,ds^2$. For the graph $r=\varepsilon a(s)$, its tangent acceleration in the $r$ direction is
$\varepsilon(a''-a)+O(\varepsilon^2)$, and its outward normal is $-\partial_r+O(\varepsilon)$. Hence
$k_g=\varepsilon(a-a'')+O(\varepsilon^2)$.
Pull the equation back to a fixed domain. Its derivative in $u$ at $(\varepsilon,u)=(0,0)$ is the Neumann operator
$h\mapsto((\Delta_g-2)h,\partial_nh)$. This is an isomorphism between the usual H\"older spaces: the form
$\int(|\nabla h|^2+2h^2)$ is coercive, so the weak solution exists uniquely, and local elliptic regularity up to a smooth boundary gives the asserted regularity. The implicit-function theorem gives smooth dependence on $\varepsilon$. Terms differentiating the pulled-back operator vanish on $u_0=0$. This proves the variation formula. It is independent of the chosen identification for the same reason.

We give a direct estimate establishing the sign in the displayed family. Let $h_0$ be the fixed solution with normal derivative $1$ on each boundary. Let $v_n$ solve $(\Delta_g-2)v_n=0$, with normal derivative $(1+n^2)\cos(ns)$ on the first boundary and zero on the other two. Then $h=h_0+v_n/2$.
Choose a fixed embedded collar $0\le r\le\delta$ of the first boundary and a smooth cutoff $\eta(r)$ equal to $1$ near zero and to $0$ near $\delta$. Define, extending by zero outside the collar,
\[
 w_n(r,s)=(n+n^{-1})\eta(r)e^{-nr}\cos(ns).
\]
It has exactly the same normal derivative as $v_n$. Put $R_n=(\Delta_g-2)w_n$. Since
\[
 \Delta_g=\partial_r^2+\tanh r\,\partial_r
                      +\operatorname{sech}^2r\,\partial_s^2,
\]
where $\eta=1$ the residual is
\[
 R_n=(n+n^{-1})e^{-nr}
       (n^2\tanh^2r-n\tanh r-2)\cos(ns).
\]
The terms involving derivatives of $\eta$ are supported away from zero and are exponentially small times a fixed polynomial in $n$. Using $|\tanh r|\le r$, changing variable $t=nr$, and using the bounded area-density factor $\cosh r$ on this collar, we obtain
\[
 \|R_n\|_{L^2(M)}^2=O(n),
 \qquad
 \int_0^\infty(t^2+t+2)^2e^{-2t}\,dt=\frac{23}{4}.
\]
All constants here depend only on the fixed pant and cutoff, not on $n$.

For clarity we justify the estimate on the correction, including its boundary. Double $M$ across its geodesic boundary. Its metric extends smoothly because $\cosh r$ is even. Extend $R_n$ evenly as an $L^2$ function on that closed hyperbolic surface. The weak solution of
$(\Delta-2)e_n=-R_n$ on the double is in $H^2$ and is even, so its restriction has homogeneous Neumann data and equals $v_n-w_n$. Integration by parts gives $\|e_n\|_2\le\|R_n\|_2/2$ and a bound for $\|\nabla e_n\|_2$ by a fixed multiple of $\|R_n\|_2$. The integrated Bochner identity, first for smooth functions and then by smooth approximation in $H^2$, on a curvature-minus-one surface is
\[
 \|\nabla^2e_n\|_2^2=\|\Delta e_n\|_2^2+\|\nabla e_n\|_2^2.
\]
Together with $\Delta e_n=2e_n-R_n$ it bounds the full $H^2$ norm by $C\|R_n\|_2$. Finally $H^2$ embeds in continuous functions on a compact surface: in each coordinate chart this follows from Fourier inversion and Cauchy--Schwarz with the integrable weight $(1+|\xi|^2)^{-2}$, and a finite partition of unity gives the global bound. Consequently
\[
 \|v_n-w_n\|_\infty=O(\sqrt n).
\]
At the first boundary point with $s=\pi/n$ this yields
\[
 h(s)=h_0(s)-\tfrac12(n+n^{-1})+O(\sqrt n)<0
\]
for every sufficiently large $n$. Fix any such $n$. Since $a_1\ge1/2$, all three collars are everywhere inward and are embedded and disjoint for small $\varepsilon$. Their boundaries are analytic. The variation formula then gives $u_\varepsilon<0$ at the corresponding boundary point for all small positive $\varepsilon$.

Finally the old area of the retained surface is strictly less than $-2\pi\chi(M)$, while the new area equals that number. Hence $\int_{M'}(e^{2u_\varepsilon}-1)dA_g>0$, which forces $u_\varepsilon>0$ somewhere. This proves that no pointwise order holds in this family, and completes the general-shape answer to 67/021.
\end{proof}
